% GENERATED FILE. Edit canonical lessons, then run npm run generate:books.
% canonical-source-hash: fnv1a64:2bb85a7138703889
% canonical-lessons: ZH-C431-jieshuo, ZH-C431-chunyou, ZH-W431-you-tour, ZH-C431-youchuan, ZH-C431-duo, ZH-C431-heying

\chapter{A Spring Outing}
\label{ch:zh-a-spring-outing}
\hlchaptermodality{431}

\begin{chapteropening}
\textbf{By the end of this chapter:} \emph{I can follow a guide's commentary, go on a spring outing, take a pleasure boat, say there are many and pose for a group photo.}

Complete the last lesson of chapter 431: I can follow a guide's commentary, go on a spring outing, take a pleasure boat, say there are many and pose for a group photo.
\end{chapteropening}

\section[\texorpdfstring{\zh{解说} (jiěshuō)}{解说 (jiěshuō)}]{\zh{解说} (jiěshuō) — to give a running commentary}

\label{lesson:ZH-C431-jieshuo}



\begin{quote}
Before the new one: say the Mandarin for to take a selfie, then the Mandarin for the peak season.
\end{quote}

\subsection*{You'll want to know: \zh{解说}}
\textbf{\zh{解说}} — \emph{jiěshuō} — \textquotedblleft{}to give a running commentary\textquotedblright{}.

\begin{quote}
\emph{dǎoyóu zhèngzài gěi wǒmen jiěshuō} — The guide is giving us a commentary.
\end{quote}

\begin{cousinweb}
\textbf{\zh{解说}} (\emph{jiěshuō}) is built from \zh{解} 'to explain' and \zh{说} 'to speak', both inherited from Middle Chinese, \emph{keaX} and \emph{sywet} (Guangyun, in Baxter's transcription), which glosses the first 'to explain' and the second 'to tell'.
\end{cousinweb}

\begin{grammarlens}[title={What you've built}]
One more word for getting around.
\end{grammarlens}

\subsection*{Your turn}
\begin{itemize}
\raggedright
  \item \emph{Say it:} \emph{jiěshuō}
  \item \emph{Say it:} \emph{jiěshuō}, once more
  \item \emph{Say it:} \emph{dǎoyóu zhèngzài gěi wǒmen jiěshuō} — The guide is giving us a commentary.
  \item \emph{Recall:} say the Mandarin for to take a selfie, then the Mandarin for the peak season, then say \emph{jiěshuō} again
\end{itemize}

\begin{tcolorbox}[breakable,colback=teal!4,colframe=teal!35!black,title={Before you move on}]
What does \zh{解说} mean? (\textquotedblleft{}To give a running commentary\textquotedblright{}.) Say it once more.
\end{tcolorbox}
\section[\texorpdfstring{\zh{春游} (chūnyóu)}{春游 (chūnyóu)}]{\zh{春游} (chūnyóu) — spring outing}

\label{lesson:ZH-C431-chunyou}



\begin{quote}
Before the new one: say the Mandarin for the peak season, then the Mandarin for to give a running commentary.
\end{quote}

\subsection*{You'll want to know: \zh{春游}}
\textbf{\zh{春游}} — \emph{chūnyóu} — \textquotedblleft{}spring outing\textquotedblright{}.

\begin{quote}
\textbf{\zh{我们坐火车去春游}} \emph{wǒmen zuò huǒchē qù chūnyóu} — We're going on a spring outing by train.
\end{quote}

The new character, in pieces: \zh{游} — twelve strokes in three pieces: on the left the water piece of \zh{洗}, three strokes; in the middle a dot, a bar, a stroke that turns down and hooks, and a stroke falling to the left; on the right a short stroke falling to the left and a bar, with \zh{子} under them. The next lesson puts it in your hand.

\begin{cousinweb}
\textbf{\zh{春游}} (\emph{chūnyóu}) is built from \zh{春} 'spring' and \zh{游} 'to roam', both inherited from Middle Chinese, \emph{tsyhwin} and \emph{yuw} (Guangyun, in Baxter's transcription), which glosses the second 'to float' and 'to let loose'; the sense 'to tour' is usually explained as a widening of it.
\end{cousinweb}

\begin{grammarlens}[title={What you've built}]
One more word for getting around.
\end{grammarlens}

\subsection*{Your turn}
\begin{itemize}
\raggedright
  \item \emph{Say it:} \emph{chūnyóu}
  \item \emph{Say it:} \emph{chūnyóu}, once more
  \item \emph{Say it:} \emph{wǒmen zuò huǒchē qù chūnyóu} — We're going on a spring outing by train.
  \item \emph{Recall:} say the Mandarin for the peak season, then the Mandarin for to give a running commentary, then say \emph{chūnyóu} again
\end{itemize}

\begin{tcolorbox}[breakable,colback=teal!4,colframe=teal!35!black,title={Before you move on}]
What does \zh{春游} mean? (\textquotedblleft{}Spring outing\textquotedblright{}.) Say it once more.
\end{tcolorbox}
\section[\texorpdfstring{\zh{游} (yóu)}{游 (yóu)}]{\zh{游} — the water piece, a middle piece, then \zh{子}}

\label{lesson:ZH-W431-you-tour}



\begin{quote}
The word first, then the character inside it.

\begin{itemize}
\raggedright
  \item \emph{Recall:} say \emph{chūnyóu} — \textbf{R1}, one lesson back, before it decays
\end{itemize}

\begin{quote}
\zh{洗} \zh{子}
\end{quote}

\emph{Read it:} each one aloud, with nothing said first — \emph{xǐ}, \emph{zǐ}
\end{quote}

\subsection*{Script — the water piece, a middle piece, then \zh{子}}
\hlblockfigure{\detokenize{figures/ZH-W431-you-tour-filmstrip.pdf}}{How \zh{游} is written, stroke by stroke}

\begin{quote}
\zh{游}
\end{quote}

\textbf{Three pieces, left to right.} Count the pieces, not the strokes.

\textbf{The left piece, the water piece, three strokes}, as in \zh{洗}: two dots, then a stroke that rises to the right.

\textbf{The middle piece, four strokes.} A dot on top. A bar under it. Under the bar, a stroke that runs a little to the right, turns down and hooks to the left at the foot. Then a stroke falling to the left, from under the bar.

\textbf{The right piece, five strokes.} A short stroke falling to the left and a bar, like a small roof. Under them \zh{子} zǐ: a bar that turns down to the left, a vertical with a hook, and a bar across.

3 + 4 + 5 = twelve.

\subsection*{Your turn}
\begin{itemize}
\raggedright
  \item \emph{Write it:} \zh{游}, all twelve, with the model in front of you
  \item \emph{Cover:} write it again from memory, and count the pieces, not the strokes
\end{itemize}

\begin{tcolorbox}[breakable,colback=teal!4,colframe=teal!35!black,title={Before you move on}]
How many pieces does \zh{游} have? (\textbf{Three: three strokes, four and five}.) Which written character sits at the bottom right? (\textbf{\zh{子}, three strokes}.)

Source: \href{https://raw.githubusercontent.com/chanind/hanzi-writer-data/68d10a4b21150cae5e1ebbd223eed289cf32d90c/data/\%E6\%B8\%B8.json}{Hanzi Writer Data \zh{游}, pinned stroke paths and medians}.
\end{tcolorbox}
\section[\texorpdfstring{\zh{游船} (yóuchuán)}{游船 (yóuchuán)}]{\zh{游船} (yóuchuán) — pleasure boat}

\label{lesson:ZH-C431-youchuan}



\begin{quote}
Before the new one: say the Mandarin for to give a running commentary, then the Mandarin for a spring outing.
\end{quote}

\subsection*{You'll want to know: \zh{游船}}
\textbf{\zh{游船}} — \emph{yóuchuán} — \textquotedblleft{}pleasure boat\textquotedblright{}.

\begin{quote}
\emph{wǒmen zuò yóuchuán qù hú duìmiàn} — We're taking a pleasure boat across the lake.
\end{quote}

\begin{cousinweb}
\textbf{\zh{游船}} (\emph{yóuchuán}) is built from \zh{游} 'to roam' and \zh{船} 'a boat', both inherited from Middle Chinese, \emph{yuw} and \emph{zywen} (Guangyun, in Baxter's transcription).
\end{cousinweb}

\begin{grammarlens}[title={What you've built}]
One more word for getting around.
\end{grammarlens}

\subsection*{Your turn}
\begin{itemize}
\raggedright
  \item \emph{Say it:} \emph{yóuchuán}
  \item \emph{Say it:} \emph{yóuchuán}, once more
  \item \emph{Say it:} \emph{wǒmen zuò yóuchuán qù hú duìmiàn} — We're taking a pleasure boat across the lake.
  \item \emph{Recall:} say the Mandarin for to give a running commentary, then the Mandarin for a spring outing, then say \emph{yóuchuán} again
\end{itemize}

\begin{tcolorbox}[breakable,colback=teal!4,colframe=teal!35!black,title={Before you move on}]
What does \zh{游船} mean? (\textquotedblleft{}Pleasure boat\textquotedblright{}.) Say it once more.
\end{tcolorbox}
\section[\texorpdfstring{\zh{多} (duō)}{多 (duō)}]{\zh{多} (duō) — many; much; a lot of}

\label{lesson:ZH-C431-duo}



\begin{quote}
Before the new one: say the Mandarin for a spring outing, then the Mandarin for a pleasure boat.
\end{quote}

\subsection*{You'll want to know: \zh{多}}
\textbf{\zh{多}} — \emph{duō} — \textquotedblleft{}many; much; a lot of\textquotedblright{}.

\begin{quote}
\textbf{\zh{春天这儿的游客很多}} \emph{chūntiān zhèr de yóukè hěn duō} — In spring there are a lot of visitors here.
\end{quote}

Before a noun it is usually \zh{很多} hěn duō: \zh{很多人} 'a lot of people'.

\begin{cousinweb}
\textbf{\zh{多}} (\emph{duō}) is inherited from Middle Chinese \emph{ta} (Guangyun, in Baxter's transcription), which glosses it 'many'.
\end{cousinweb}

\begin{grammarlens}[title={What you've built}]
One more word for getting around.
\end{grammarlens}

\subsection*{Your turn}
\begin{itemize}
\raggedright
  \item \emph{Say it:} \emph{duō}
  \item \emph{Say it:} \emph{duō}, once more
  \item \emph{Say it:} \emph{chūntiān zhèr de yóukè hěn duō} — In spring there are a lot of visitors here.
  \item \emph{Recall:} say the Mandarin for a spring outing, then the Mandarin for a pleasure boat, then say \emph{duō} again
\end{itemize}

\begin{tcolorbox}[breakable,colback=teal!4,colframe=teal!35!black,title={Before you move on}]
What does \zh{多} mean? (\textquotedblleft{}Many; much; a lot of\textquotedblright{}.) Say it once more.
\end{tcolorbox}
\section[héyǐng]{héyǐng — to have a group photo taken; to pose for a group photo}

\label{lesson:ZH-C431-heying}



\begin{quote}
Before the new one: say the Mandarin for a pleasure boat, then the Mandarin for many.
\end{quote}

\subsection*{You'll want to know: héyǐng}
\emph{héyǐng} — \textquotedblleft{}to have a group photo taken; to pose for a group photo\textquotedblright{}.

\begin{quote}
\emph{wǒmen zài jiǔdiàn ménkǒu héyǐng} — We're having a group photo taken outside the hotel.
\end{quote}

Its characters come later; for now it is a word for the ear.

\begin{cousinweb}
\emph{héyǐng} is built from \emph{hé} 'together' and \emph{yǐng} 'a likeness', both inherited from Middle Chinese, \emph{hop} and \emph{'jaengX} (Guangyun, in Baxter's transcription), which glosses the first 'to join' and the second as a shape and its shadow; the sense 'a photograph' is modern.
\end{cousinweb}

\begin{grammarlens}[title={What you've built}]
One more word for getting around.
\end{grammarlens}

\subsection*{Your turn}
\begin{itemize}
\raggedright
  \item \emph{Say it:} \emph{héyǐng}
  \item \emph{Say it:} \emph{héyǐng}, once more
  \item \emph{Say it:} \emph{wǒmen zài jiǔdiàn ménkǒu héyǐng} — We're having a group photo taken outside the hotel.
  \item \emph{Recall:} say the Mandarin for a pleasure boat, then the Mandarin for many, then say \emph{héyǐng} again
\end{itemize}

\begin{tcolorbox}[breakable,colback=teal!4,colframe=teal!35!black,title={Before you move on}]
What does héyǐng mean? (\textquotedblleft{}To have a group photo taken; to pose for a group photo\textquotedblright{}.) Say it once more.
\end{tcolorbox}
